\documentclass{article}

\usepackage[T2A]{fontenc}
\usepackage[utf8]{inputenc}
\usepackage{helvet}
\usepackage{geometry}
\usepackage{xcolor}
\usepackage{eso-pic}
\usepackage{fancyhdr}
\usepackage{parskip}
\usepackage{microtype}
\usepackage[english, ukrainian]{babel}
\usepackage{url}
\usepackage{amsmath}
\usepackage{amssymb}
\usepackage{amsthm}
\usepackage{longtable}
\usepackage{booktabs}
\usepackage{hyperref}
\usepackage{titlesec}

\definecolor{nato}{HTML}{293757}
\definecolor{footergray}{gray}{0.65}

\geometry{a4paper,left=3cm,right=3cm,top=3cm,bottom=3cm}
\renewcommand{\familydefault}{\sfdefault}
\setcounter{secnumdepth}{3}

\titleformat{\section}
  {\color{nato}\Huge\bfseries}{\thesection.}{1em}{}
\titlespacing*{\section}{0pt}{30pt}{12pt}

\titleformat{\subsection}
  {\color{nato}\LARGE\bfseries}{\thesubsection.}{1em}{}
\titlespacing*{\subsection}{0pt}{20pt}{8pt}

\titleformat{\subsubsection}
  {\color{nato}\Large\bfseries}{\thesubsubsection.}{1em}{}
\titlespacing*{\subsubsection}{0pt}{10pt}{6pt}

\fancyhf{}
\fancyhead[R]{\small\color{footergray} 9 червня 2026}
\fancyfoot[L]{\small\color{footergray} Програма модернізації ЄСІКС. Цілі та Ключові Результати (OKR)}
\fancyfoot[R]{\small\color{footergray}\thepage}

\newcommand\HUGE[1]{\fontsize{40}{40}\selectfont #1}

\renewcommand{\headrulewidth}{0pt}
\renewcommand{\footrulewidth}{0.4pt}
\renewcommand{\footrule}{\color{footergray}\hrule width \textwidth height 0.4pt}

\begin{document}

\pagestyle{fancy}

\begin{titlepage}
  \AddToShipoutPictureBG*{\AtPageLowerLeft{\color{nato}\rule{\paperwidth}{\paperheight}}}
  \color{white}\thispagestyle{empty}\vspace*{3cm}\centering
  {\Huge \textbf{ЦІЛІ ТА КЛЮЧОВІ РЕЗУЛЬТАТИ} \par} \vspace{1cm}
  {\HUGE \textbf{Програма модернізації ЄСІКС} \par} \vspace{2cm}
  {\LARGE \textbf{Узгоджений стратегічний план \\ розробки та впровадження підсистем ЄСІКС} \par} \vspace{1.5cm}
  {\large\textbf{Цикл планування: Q3--Q4 2026} \par}
  \vfill
  {\large 2026 \copyright\ ДП «Інформаційні судові системи» \par}
\end{titlepage}

\newpage
\thispagestyle{empty}
\begin{center}
  {\Large\color{nato}\bfseries ЛИСТ ЗАТВЕРДЖЕННЯ ТА КОНТРОЛЮ ВЕРСІЙ\par}
\end{center}
\vspace{20pt}

\begin{longtable}{@{}p{4.0cm}p{10.2cm}@{}}
\toprule
\textbf{Реквізит документа} & \textbf{Значення та відомості} \\
\midrule
\textbf{Назва документа} & Цілі та Ключові Результати (OKR) модернізації ЄСІКС \\
\textbf{Об'єкт планування} & Єдина судова інформаційно-телекомунікаційна система (ЄСІКС / ЄСІКС) \\
\textbf{Версія документа} & v1.0.0 (Офіційна редакція) \\
\textbf{Дата затвердження} & 09.06.2026 \\
\textbf{Цільовий період} & III--IV квартали 2026 року \\
\textbf{Статус документа} & Затверджено робочою групою з модернізації \\
\textbf{Відповідність стандартам} & ISO 21502:2020 (Управління проектами), ISO 9001:2015 (Якість), ISO 31000 (Ризики), ISO 15489 (Документування) \\
\textbf{Розробники} & Робоча група з питань модернізації та цифровізації судочинства \\
\textbf{Погоджено} & Архітектурний комітет ЄСІКС, Керівництво програми модернізації \\
\textbf{Затвердив} & Генеральний директор ДП «Інформаційні судові системи» \\
\bottomrule
\end{longtable}

\vspace{20pt}
\begin{center}
  {\large\color{nato}\bfseries Історія змін\par}
\end{center}
\vspace{10pt}
\begin{longtable}{@{}p{1.5cm}p{2.0cm}p{3.0cm}p{7.0cm}@{}}
\toprule
\textbf{Версія} & \textbf{Дата} & \textbf{Автор} & \textbf{Опис внесених змін} \\
\midrule
0.1.0 & 25.05.2026 & Робоча група & Первинний проєкт цілей та ключових результатів \\
0.9.0 & 02.06.2026 & Архітектор ЄСІКС & Додано архітектурні орієнтири та карти інтеграцій \\
1.0.0 & 09.06.2026 & Керівник програми & Затверджено офіційну версію для виконання \\
\bottomrule
\end{longtable}

\newpage
\begin{center}
  {\Huge\color{nato}\bfseries Зміст\par}
\end{center}
\vspace{40pt}
\tableofcontents

\newpage

\section{Вступ та загальні положення}

\subsection{Контекст та передумови}
Цей документ визначає Цілі та Ключові Результати (Objectives and Key Results, OKR) для програми модернізації Єдиної судової інформаційно-телекомунікаційної системи (ЄСІКС / ЄСІКС) на період III--IV кварталів 2026 року. Модернізація ЄСІКС є критично важливим державним проектом, спрямованим на підвищення доступності, прозорості та ефективності правосуддя в Україні за допомогою сучасних цифрових рішень.

Використання методології OKR дозволяє сфокусувати зусилля розробників, аналітиків та керівників на пріоритетних напрямах, забезпечити вимірюваність результатів та гнучкість управління в умовах високої невизначеності та технологічних змін.

\subsection{Відповідність міжнародним стандартам управління}
Хоча прямого міжнародного стандарту ISO для OKR-документів не існує, даний документ структуровано з урахуванням вимог та найкращих практик наступних міжнародних і національних стандартів:
\begin{itemize}
    \item \textbf{ISO 21502:2020} (\emph{Project, programme and portfolio management --- Guidance on project management}): стандартизує постановку цілей, моніторинг прогресу, визначення залежностей та очікуваних результатів.
    \item \textbf{ISO 9001:2015} (\emph{Quality management systems}): вимагає чіткості формулювання цілей якості та їх вимірюваності на кожному рівні організації.
    \item \textbf{ISO 31000} (\emph{Risk management}): регулює виявлення, оцінку та управління ризиками при досягненні поставлених результатів (див. Додаток А).
    \item \textbf{ISO 15489-1} (\emph{Information and documentation --- Records management}): визначає правила структурування та контролю версійності службових документів.
\end{itemize}

\subsection{Структура документа}
Відповідно до рекомендацій OKR Institute та стандартів управління проектами, кожна ціль (Objective) містить перелік вимірюваних ключових результатів (Key Results). Кожен ключовий результат має:
\begin{itemize}
    \item \textbf{Опис очікуваного стану} (конкретні бізнес- або технічні критерії);
    \item \textbf{Власника} (відповідальну посадову особу або роль);
    \item \textbf{Показник прогресу / Метрику} (вимірюється від 0\% до 100\%);
    \item \textbf{Граничний термін} досягнення цільового значення;
    \item \textbf{Критерії виконання} (складові роботи та підтверджуючі артефакти).
\end{itemize}

\newpage

\section{Цілі та Ключові Результати (OKR)}

\subsection{Ціль 1 (O1)}

\textbf{Формулювання цілі:} Перетворити бачення ЄСІКСу на керовану основу, з якої команда може впевнено стартувати модернізацію.

\textbf{Контекст та обґрунтування:}
Дана ціль спрямована на усунення невизначеності щодо складу підсистем, визначення їхніх взаємозв'язків, версійності та чітких критеріїв запуску. Це створює міцну управлінську та методологічну базу для всіх наступних етапів розробки та впровадження.

\subsubsection{Ключовий результат 1.1 (KR 1.1)}

\textbf{Формулювання результату:} З 0 до 100\% підсистем ЄСІКС мають погоджену карту залежностей, версійність, критичні передумови та статус готовності до запуску.

\noindent\begin{tabular}{@{}p{2.2cm}p{4.5cm}p{2.2cm}p{4.5cm}@{}}
\toprule
\textbf{Власник:} & Архітектор ЄСІКС & \textbf{Метрика:} & \% підсистем з картою \\
\textbf{Термін:} & 30 вересня 2026 р. & \textbf{Статус:} & В процесі \\
\bottomrule
\end{tabular}

\vspace{10pt}
\textbf{Складові роботи та критерії виконання:}
\begin{itemize}
    \item Визначено межі та призначення кожної підсистеми;
    \item Погоджено склад функціональних компонентів (підсистеми, модулі, сервіси);
    \item Визначено критичні та некритичні залежності між компонентами;
    \item Визначено вплив залежностей на можливість реалізації та запуску;
    \item Погоджено цільову послідовність реалізації та взаємозалежність етапів;
    \item Зафіксовано обмеження, передумови та блокуючі фактори;
    \item Визначено інтеграційний контур та точки взаємодії між системами.
\end{itemize}

\subsubsection{Ключовий результат 1.2 (KR 1.2)}

\textbf{Формулювання результату:} З 0 до 100\% підсистем із погодженим обсягом реалізації та моделлю поетапного розвитку.

\noindent\begin{tabular}{@{}p{2.2cm}p{4.5cm}p{2.2cm}p{4.5cm}@{}}
\toprule
\textbf{Власник:} & Власник продукту (PO) & \textbf{Метрика:} & \% підсистем з MVP \\
\textbf{Термін:} & 15 жовтня 2026 р. & \textbf{Статус:} & В процесі \\
\bottomrule
\end{tabular}

\vspace{10pt}
\textbf{Складові роботи та критерії виконання:}
\begin{itemize}
    \item Визначено мінімально достатній функціональний обсяг першого запуску (MVP);
    \item Визначено функціональність, перенесену на наступні етапи;
    \item Погоджено модель версійності та правила переходу між етапами.
\end{itemize}

\subsubsection{Ключовий результат 1.3 (KR 1.3)}

\textbf{Формулювання результату:} З 0 до 100\% підсистем із погодженими критеріями готовності до запуску по всіх вимірах.

\noindent\begin{tabular}{@{}p{2.2cm}p{4.5cm}p{2.2cm}p{4.5cm}@{}}
\toprule
\textbf{Власник:} & QA Lead / Юрист & \textbf{Метрика:} & \% підсистем з критеріями \\
\textbf{Термін:} & 31 жовтня 2026 р. & \textbf{Статус:} & Не розпочато \\
\bottomrule
\end{tabular}

\vspace{10pt}
\textbf{Складові роботи та критерії виконання:}
\begin{itemize}
    \item Визначено юридичні передумови запуску (нормативні зміни, умови введення в експлуатацію);
    \item Визначено технічні передумови запуску;
    \item Визначено вимоги до операційної готовності;
    \item Погоджено критерії переходу до запуску, пілотування або наступного етапу;
    \item Зафіксовано ризики, що можуть вплинути на готовність.
\end{itemize}

\subsubsection{Ключовий результат 1.4 (KR 1.4)}

\textbf{Формулювання результату:} З 0 до 100\% підсистем із зафіксованим поточним статусом, цільовим станом та наступним управлінським кроком.

\noindent\begin{tabular}{@{}p{2.2cm}p{4.5cm}p{2.2cm}p{4.5cm}@{}}
\toprule
\textbf{Власник:} & Керівник програми ЄСІКС & \textbf{Метрика:} & \% підсистем з планом кроків \\
\textbf{Термін:} & 15 листопада 2026 р. & \textbf{Статус:} & Не розпочато \\
\bottomrule
\end{tabular}

\vspace{10pt}
\textbf{Складові роботи та критерії виконання:}
\begin{itemize}
    \item Визначено поточний статус реалізації;
    \item Визначено цільовий статус на кінець OKR-циклу;
    \item Погоджено умови переходу між статусами;
    \item Визначено відповідального за просування підсистеми;
    \item Погоджено орієнтовний горизонт запуску або переходу до наступного етапу;
    \item Зафіксовано рішення, яке потрібно ухвалити для подальшого руху.
\end{itemize}

\newpage

\subsection{Ціль 2 (O2)}

\textbf{Формулювання цілі:} Сформувати цільову архітектурну модель ЄСІКС та забезпечити керований перехід до її реалізації.

\textbf{Контекст та обґрунтування:}
Ця ціль спрямована на вибудовування стійкої та масштабованої системної архітектури, регламентацію міжсистемної взаємодії та безпечну міграцію зі застарілих рішень на цільові платформи.

\subsubsection{Ключовий результат 2.1 (KR 2.1)}

\textbf{Формулювання результату:} З 0 до 100\% підсистем із погодженою цільовою архітектурною моделлю та правилами взаємодії.

\noindent\begin{tabular}{@{}p{2.2cm}p{4.5cm}p{2.2cm}p{4.5cm}@{}}
\toprule
\textbf{Власник:} & Провідний архітектор & \textbf{Метрика:} & \% підсистем з арх. схемою \\
\textbf{Термін:} & 31 жовтня 2026 р. & \textbf{Статус:} & В процесі \\
\bottomrule
\end{tabular}

\vspace{10pt}
\textbf{Складові роботи та критерії виконання:}
\begin{itemize}
    \item Визначено архітектурні принципи побудови та взаємодії;
    \item Визначено правила міжсистемної взаємодії та обміну даними.
\end{itemize}

\subsubsection{Ключовий результат 2.2 (KR 2.2)}

\textbf{Формулювання результату:} З 0 до 100\% ключових інтеграцій із погодженими правилами реалізації та відповідальністю сторін.

\noindent\begin{tabular}{@{}p{2.2cm}p{4.5cm}p{2.2cm}p{4.5cm}@{}}
\toprule
\textbf{Власник:} & Інтеграційний архітектор & \textbf{Метрика:} & \% погоджених інтеграцій \\
\textbf{Термін:} & 30 листопада 2026 р. & \textbf{Статус:} & В процесі \\
\bottomrule
\end{tabular}

\vspace{10pt}
\textbf{Складові роботи та критерії виконання:}
\begin{itemize}
    \item Визначено перелік критичних інтеграцій;
    \item Визначено власників інтеграцій та відповідальні сторони;
    \item Визначено мінімальний інтеграційний контур для запуску;
    \item Погоджено етапність реалізації інтеграцій.
\end{itemize}

\subsubsection{Ключовий результат 2.3 (KR 2.3)}

\textbf{Формулювання результату:} З 0 до 100\% підсистем із погодженим сценарієм переходу від поточного до цільового стану.

\noindent\begin{tabular}{@{}p{2.2cm}p{4.5cm}p{2.2cm}p{4.5cm}@{}}
\toprule
\textbf{Власник:} & PM / Керівник програми & \textbf{Метрика:} & \% підсистем зі сценаріями \\
\textbf{Термін:} & 15 грудня 2026 р. & \textbf{Статус:} & Не розпочато \\
\bottomrule
\end{tabular}

\vspace{10pt}
\textbf{Складові роботи та критерії виконання:}
\begin{itemize}
    \item Визначено підсистеми, що залишаються чинними на перехідному етапі;
    \item Визначено підсистеми, що підлягають модернізації або заміні;
    \item Визначено правила співіснування старих і нових компонентів;
    \item Визначено умови відключення застарілих компонентів;
    \item Погоджено механізми міграції функціональності та даних;
    \item Визначено критерії завершення переходу до цільового стану.
\end{itemize}

\subsubsection{Ключовий результат 2.4 (KR 2.4)}

\textbf{Формулювання результату:} З 0 до 100\% підсистем із погодженими профілями безпеки та порядком впровадження контролів відповідно до NIST SP 800-53 / НД ТЗІ 3.6-006-24.

\noindent\begin{tabular}{@{}p{2.2cm}p{4.5cm}p{2.2cm}p{4.5cm}@{}}
\toprule
\textbf{Власник:} & CTO / QA Lead & \textbf{Метрика:} & \% підсистем з профілем \\
\textbf{Термін:} & 31 грудня 2026 р. & \textbf{Статус:} & Не розпочато \\
\bottomrule
\end{tabular}

\vspace{10pt}
\textbf{Складові роботи та критерії виконання:}
\begin{itemize}
    \item Проведено категоризацію підсистем відповідно до FIPS 199 / FIPS 200 та визначено базовий набір контролів безпеки (Baseline);
    \item Розроблено та погоджено Цільові профілі безпеки підсистем на основі Галузевого профілю судової системи (Court Industry Profile) згідно НД ТЗІ 3.6-006-24;
    \item Визначено вимоги до імплементації Опорного монітора (Reference Monitor, AC-25) для забезпечення неможливості обходу політик доступу (non-bypassable);
    \item Розроблено плани забезпечення безпеки систем (System Security Plan, SSP) відповідно до NIST SP 800-37 (Risk Management Framework);
    \item Погоджено двоетапну процедуру розділення обов'язків (Segregation of Duties) для створення та атестації КСЗІ з двома незалежними провайдерами-ліцензіатами;
    \item Розроблено Програми державної експертизи КСЗІ для проведення незалежного технічного аудиту та інструментального тестування (NIST SP 800-53A);
    \item Визначено критерії готовності до отримання Атестата відповідності від Держспецзв'язку (або надання Authority to Operate - ATO).
\end{itemize}

\newpage

\subsection{Ціль 3 (O3)}

\textbf{Формулювання цілі:} Підготувати організацію до впровадження OKR-підходу в межах ЄСІКС.

\textbf{Контекст та обґрунтування:}
Ця ціль спрямована на організаційну трансформацію та навчання команд, встановлення ритму синхронізації та забезпечення прозорості у відслідковуванні досягнення цілей.

\subsubsection{Ключовий результат 3.1 (KR 3.1)}

\textbf{Формулювання результату:} З 0 до 100\% ключових учасників програми розуміють логіку OKR, ролі, цикл планування та правила формулювання результатів.

\noindent\begin{tabular}{@{}p{2.2cm}p{4.5cm}p{2.2cm}p{4.5cm}@{}}
\toprule
\textbf{Власник:} & PMO Lead / Agile Coach & \textbf{Метрика:} & \% сертифікованих учасників \\
\textbf{Термін:} & 31 серпня 2026 р. & \textbf{Статус:} & В процесі \\
\bottomrule
\end{tabular}

\vspace{10pt}
\textbf{Складові роботи та критерії виконання:}
\begin{itemize}
    \item Проведено вступне пояснення логіки OKR для команди;
    \item Узгоджено різницю між Objective, Key Result, ініціативами, задачами та roadmap;
    \item Пояснено ролі учасників OKR-процесу;
    \item Погоджено цикл планування, перегляду та оновлення OKR;
    \item Визначено правила формулювання якісних Objective та вимірюваних KR;
    \item Зафіксовано приклади правильних і неправильних формулювань для ЄСІКС;
    \item Погоджено єдиний шаблон опису OKR.
\end{itemize}

\subsubsection{Ключовий результат 3.2 (KR 3.2)}

\textbf{Формулювання результату:} З 0 до 100\% продуктових напрямів мають визначених відповідальних, ритм синхронізації та формат звітності.

\noindent\begin{tabular}{@{}p{2.2cm}p{4.5cm}p{2.2cm}p{4.5cm}@{}}
\toprule
\textbf{Власник:} & PMO Lead / Scrum Master & \textbf{Метрика:} & \% підсистем з процесом OKR \\
\textbf{Термін:} & 15 вересня 2026 р. & \textbf{Статус:} & В процесі \\
\bottomrule
\end{tabular}

\vspace{10pt}
\textbf{Складові роботи та критерії виконання:}
\begin{itemize}
    \item Визначено перелік продуктових напрямів / підсистем, які входять в OKR-контур;
    \item Визначено відповідальних за кожен продуктовий напрям;
    \item Визначено відповідальних за оновлення статусів і прогресу;
    \item Погоджено формат статусної звітності;
    \item Визначено правила ескалації ризиків, блокерів і рішень;
    \item Погоджено єдине место ведення OKR-статусів.
\end{itemize}

\newpage

\subsection{Ціль 4 (O4)}

\textbf{Формулювання цілі:} Підготувати ЄСІКС до контрольованого пілотування та прийняття рішень щодо масштабування перших рішень.

\textbf{Контекст та обґрунтування:}
Ця ціль спрямована на проведення тестових випробувань модернізованих підсистем у реальних умовах, мінімізацію технічних та юридичних ризиків перед повноцінним запуском.

\subsubsection{Ключовий результат 4.1 (KR 4.1)}

\textbf{Формулювання результату:} З 0 до 100\% підсистем / сервісів першої хвилі пілотування мають погоджений контур пілоту.

\noindent\begin{tabular}{@{}p{2.2cm}p{4.5cm}p{2.2cm}p{4.5cm}@{}}
\toprule
\textbf{Власник:} & Керівник впровадження & \textbf{Метрика:} & \% сервісів з контуром пілоту \\
\textbf{Термін:} & 15 листопада 2026 р. & \textbf{Статус:} & Не розпочато \\
\bottomrule
\end{tabular}

\vspace{10pt}
\textbf{Складові роботи та критерії виконання:}
\begin{itemize}
    \item Визначено перелік підсистем і сервісів, придатних до пілотування;
    \item Визначено межі пілоту та обсяг функціональності;
    \item Визначено перелік користувачів та ролей;
    \item Визначено бізнес-процеси, які перевіряються під час пілоту;
    \item Визначено залежності та обмеження проведення пілоту;
    \item Визначено очікувану цінність і цільовий результат пілотування.
\end{itemize}

\subsubsection{Ключовий результат 4.2 (KR 4.2)}

\textbf{Формулювання результату:} З 0 до 100\% пілотів мають погоджені умови готовності до запуску.

\noindent\begin{tabular}{@{}p{2.2cm}p{4.5cm}p{2.2cm}p{4.5cm}@{}}
\toprule
\textbf{Власник:} & CTO / QA Lead & \textbf{Метрика:} & \% пілотів з критеріями запуску \\
\textbf{Термін:} & 30 листопада 2026 р. & \textbf{Статус:} & Не розпочато \\
\bottomrule
\end{tabular}

\vspace{10pt}
\textbf{Складові роботи та критерії виконання:}
\begin{itemize}
    \item Визначено функціональні критерії готовності;
    \item Визначено технічні критерії готовності;
    \item Визначено юридичні умови проведення пілоту;
    \item Визначено вимоги до операційної готовності;
    \item Визначено вимоги до безпеки та захисту даних;
    \item Визначено критерії допуску до старту пілоту;
    \item Зафіксовано ризики та механізми їх контролю.
\end{itemize}

\subsubsection{Ключовий результат 4.3 (KR 4.3)}

\textbf{Формулювання результату:} З 0 до 100\% пілотів мають погоджену модель проведення та оцінки результатів.

\noindent\begin{tabular}{@{}p{2.2cm}p{4.5cm}p{2.2cm}p{4.5cm}@{}}
\toprule
\textbf{Власник:} & Керівник впровадження & \textbf{Метрика:} & \% пілотів з моделлю оцінки \\
\textbf{Термін:} & 15 грудня 2026 р. & \textbf{Статус:} & Не розпочато \\
\bottomrule
\end{tabular}

\vspace{10pt}
\textbf{Складові роботи та критерії виконання:}
\begin{itemize}
    \item Визначено пілотні майданчики та критерії їх вибору;
    \item Визначено ролі учасників та модель підтримки;
    \item Погоджено сценарії тестування та порядок проведення;
    \item Визначено очікувані результати та критерії успіху;
    \item Визначено показники оцінки результатів;
    \item Визначено механізми збору та аналізу зворотного зв’язку;
    \item Погоджено формат ухвалення рішень за результатами пілоту.
\end{itemize}

\subsubsection{Ключовий результат 4.4 (KR 4.4)}

\textbf{Формулювання результату:} З 0 до 100\% пілотів мають погоджений сценарій переходу до масштабування або перегляду рішення.

\noindent\begin{tabular}{@{}p{2.2cm}p{4.5cm}p{2.2cm}p{4.5cm}@{}}
\toprule
\textbf{Власник:} & Керівник програми ЄСІКС & \textbf{Метрика:} & \% пілотів зі сценарієм масштабування \\
\textbf{Термін:} & 31 грудня 2026 р. & \textbf{Статус:} & Не розпочато \\
\bottomrule
\end{tabular}

\vspace{10pt}
\textbf{Складові роботи та критерії виконання:}
\begin{itemize}
    \item Визначено критерії успішності пілоту;
    \item Визначено умови переходу до наступного етапу;
    \item Визначено ризики масштабування;
    \item Визначено перелік необхідних доопрацювань;
    \item Визначено цільовий сценарій подальших дій (масштабування / повторне пілотування / доопрацювання);
    \item Погоджено наступний управлінський крок після завершення пілоту.
\end{itemize}

\newpage

\subsection{Ціль 5 (O5)}

\textbf{Формулювання цілі:} Модернізувати та стабілізувати підсистему відеоконференцзв’язку (ВКЗ) судової системи відповідно до міжнародних стандартів WebRTC (RFC 8825, 8829, 8835), медіакодеків (ITU-T H.264 / H.265) та впровадити вбудовану систему WebSocket-телеметрії.

\textbf{Контекст та обґрунтування:}
Модернізація ВКЗ спрямована на усунення ризиків нестабільної роботи підсистеми під час пікових навантажень, забезпечення сумісності медіа-трактів та мультиплексорів із галузевими стандартами, а також організацію проактивного моніторингу клієнтів і серверів за допомогою низьколатентної телеметрії на вебсокетах (RFC 6455).

\subsubsection{Ключовий результат 5.1 (KR 5.1)}

\textbf{Формулювання результату:} Переведення 100\% медіа-серверів на OpenVidu (OV) та впровадження нового типу конференції «без учасників» (без мікшування відеопотоків) для оптимізації CPU відповідно до стандартів RTP/RTCP (RFC 3550).

\noindent\begin{tabular}{@{}p{2.2cm}p{4.5cm}p{2.2cm}p{4.5cm}@{}}
\toprule
\textbf{Власник:} & Керівник напрямку ВКЗ & \textbf{Метрика:} & \% залів з OV та новими типами \\
\textbf{Термін:} & 31 грудня 2026 р. & \textbf{Статус:} & В процесі \\
\bottomrule
\end{tabular}

\vspace{10pt}
\textbf{Складові роботи та критерії виконання:}
\begin{itemize}
    \item Оновлення ПЗ OpenVidu до версії 3.2.0;
    \item Реалізація конференції типу «без учасника» з безпосереднім записом одного відеопотоку;
    \item Зменшення часу видалення кімнат конференцій після завершення засідання до 10 секунд;
    \item Підвищення ліміту одночасних записів на медіа-серверах з 10 до 18.
\end{itemize}

\subsubsection{Ключовий результат 5.2 (KR 5.2)}

\textbf{Формулювання результату:} Реалізація роботи сигнального серверу в декілька інстансів із синхронізацією через Redis Pub/Sub, відкриттям портів на TURN нодах та впровадженням вбудованої WebSocket-телеметрії (RFC 6455).

\noindent\begin{tabular}{@{}p{2.2cm}p{4.5cm}p{2.2cm}p{4.5cm}@{}}
\toprule
\textbf{Власник:} & Архітектор мережі / CTO & \textbf{Метрика:} & \% інтегрованих інстансів та покриття телеметрією \\
\textbf{Термін:} & 31 жовтня 2026 р. & \textbf{Статус:} & В процесі \\
\bottomrule
\end{tabular}

\vspace{10pt}
\textbf{Складові роботи та критерії виконання:}
\begin{itemize}
    \item Оновлення Socket.io до версії 4.8.3 на сигнальних серверах;
    \item Налаштування Redis Pub/Sub для горизонтального масштабування сигнальних серверів;
    \item Конфігурація пулу з 10 TURN нод (turn.court.gov.ua, turn1-10.court.gov.ua) з динамічним відкриттям портів WebRTC (RFC 5245, 8489, 8656);
    \item Впровадження вбудованої WebSocket-телеметрії (WSS, RFC 6455) для передачі клієнтських WebRTC Statistics (ICE state, packet loss, jitter, RTT, CPU load) у реальному часі.
\end{itemize}

\subsubsection{Ключовий результат 5.3 (KR 5.3)}

\textbf{Формулювання результату:} Розробка та тестування важкого Flutter-клієнта для автономного офлайн запису засідань з підтримкою плат BlackMagic/OBS та локальної БД.

\noindent\begin{tabular}{@{}p{2.2cm}p{4.5cm}p{2.2cm}p{4.5cm}@{}}
\toprule
\textbf{Власник:} & Провідний Flutter-розробник & \textbf{Метрика:} & \% готовності Flutter-клієнта \\
\textbf{Термін:} & 30 листопада 2026 р. & \textbf{Статус:} & В процесі \\
\bottomrule
\end{tabular}

\vspace{10pt}
\textbf{Складові роботи та критерії виконання:}
\begin{itemize}
    \item Реалізація офлайн запису відеофонограми без підключення до мережі Інтернет;
    \item Впровадження локальної БД (SQLite/LSM) для збереження дій та тимчасових протоколів;
    \item Інтеграція з платами BlackMagic та OBS Studio;
    \item Автоматична синхронізація та вивантаження локальних архівів після відновлення мережевого зв'язку.
\end{itemize}

\subsubsection{Ключовий результат 5.4 (KR 5.4)}

\textbf{Формулювання результату:} Перенесення 100\% файлів запису засідань на сховище Scality з прямим стрімінгом через Nginx та автоочищенням історичних дій.

\noindent\begin{tabular}{@{}p{2.2cm}p{4.5cm}p{2.2cm}p{4.5cm}@{}}
\toprule
\textbf{Власник:} & Адміністратор сховищ / DBA & \textbf{Метрика:} & \% інтегрованих архівів у Scality \\
\textbf{Термін:} & 15 грудня 2026 р. & \textbf{Статус:} & В процесі \\
\bottomrule
\end{tabular}

\vspace{10pt}
\textbf{Складові роботи та критерії виконання:}
\begin{itemize}
    \item Налаштування автоматичного перенесення MP4 файлів засідань на Scality;
    \item Впровадження онлайн перегляду записів через Nginx без використання PHP-серверів;
    \item Створення standby копії БД PostgreSQL 16 на РЦОД та автоматичний розподіл запитів «тільки на читання»;
    \item Впровадження очищення таблиці дій підписаних протоколів старше 3 місяців для оновлення та оптимізації БД.
\end{itemize}

\newpage

\section{Стратегічні ініціативи та завдання}
Для досягнення встановлених OKR розроблено перелік ключових стратегічних ініціатив (Initiatives) згідно з методологією ISO 21502.

\subsection{Ініціатива І--1: Аудит та картографування архітектури ЄСІКС}
\textbf{Мета:} Створення єдиної карти залежностей та інтеграцій між усіма 10-ма доменами судових та системних компонент ЄСІКС (відповідно до таксономії Політик).
\begin{itemize}
    \item Залучені підрозділи: Відділ архітектури, Аналітичний відділ, Зовнішні консультанти.
    \item Ключові завдання: Інвентаризація API-ендпоінтів, моделювання потоків даних.
    \item Строки: Червень -- Вересень 2026.
\end{itemize}

\subsection{Ініціатива І--2: Розробка нормативного та юридичного базису}
\textbf{Мета:} Підготовка законодавчих та адміністративних змін для легалізації пілотних проектів судочинства.
\begin{itemize}
    \item Залучені підрозділи: Юридичний департамент ДП «ІСС», Державна судова адміністрація (ДСА).
    \item Ключові завдання: Проекти наказів ДСА про проведення пілотування модернізованого ЄСІКС.
    \item Строки: Серпень -- Листопад 2026.
\end{itemize}

\subsection{Ініціатива І--3: Впровадження OKR та Agile-процесів}
\textbf{Мета:} Підвищення організаційної зрілості команд розробки та впровадження.
\begin{itemize}
    \item Залучені підрозділи: PMO, Scrum-майстри, продуктові команди.
    \item Ключові завдання: Проведення навчальних сесій, інтеграція OKR у JIRA та Confluence.
    \item Строки: Червень -- Серпень 2026.
\end{itemize}

\newpage

\section{Механізм управління та моніторингу OKR}
Відповідно до вимог ISO 9001:2015 щодо постійного моніторингу та оцінки досягнення цілей, в рамках програми впроваджується наступний регламент:

\subsection{Цикл синхронізації та оновлення}
Моніторинг прогресу здійснюється на трьох рівнях:
\begin{enumerate}
    \item \textbf{Щотижнева синхронізація (Weekly Check-in):} Короткі зустрічі команд (до 15 хв.) під керівництвом Scrum-майстра/Власника продукту для оновлення поточних статусів завдань та виявлення блокерів.
    \item \textbf{Щомісячний огляд прогресу (Monthly Review):} Зустріч власників KR із керівником програми ЄСІКС для оцінки прогресу (у відсотках 0--100\%) та коригування ініціатив.
    \item \textbf{Квартальна оцінка та планування (Quarterly Review \& Planning):} Фіксація фінальних оцінок за попередній цикл, ретроспектива та затвердження OKR на наступний період.
\end{enumerate}

\subsection{Система оцінювання (Scoring Guide)}
Кожен Key Result оцінюється в кінці циклу за шкалою від 0.0 до 1.0 відповідно до кращих світових практик:
\begin{longtable}{@{}p{2.0cm}p{2.0cm}p{9.5cm}@{}}
\toprule
\textbf{Оцінка} & \textbf{Колір} & \textbf{Значення та інтерпретація результату} \\
\midrule
\textbf{0.7 -- 1.0} & Зелений & \textbf{Успішно виконано.} Ціль досягнута або максимально наближена до досягнення. \\
\midrule
\textbf{0.4 -- 0.6} & Жовтий & \textbf{Частковий прогрес.} Зроблено значний поступ, але ціль не досягнута повністю через певні ризики чи блокери. \\
\midrule
\textbf{0.0 -- 0.3} & Червоний & \textbf{Незадовільний прогрес.} Жодного суттєвого просування. Потребує негайного перегляду підходу або зміни пріоритетів. \\
\bottomrule
\end{longtable}

Середній бал по кожній цілі визначається як середнє арифметичне оцінок відповідних Ключових Результатів. Загальний показник програми вище 0.7 вважається відмінним результатом.

\newpage

\section{Додатки}

\subsection{Додаток А: Матриця ризиків (ISO 31000)}
Відповідно до вимог стандарту ISO 31000, нижче наведена матриця ризиків, що можуть безпосередньо вплинути на виконання OKR програми модернізації ЄСІКС.

\begin{longtable}{@{}p{4.0cm}p{1.5cm}p{1.5cm}p{6.5cm}@{}}
\toprule
\textbf{Опис ризику} & \textbf{Ймовірн.} & \textbf{Вплив} & \textbf{Заходи з мінімізації ризику} \\
\midrule
\textbf{Юридичне затягування:} Затримка погодження нормативних змін ДСА або ВРП. & Середня & Високий & Раннє залучення юристів, розробка Тимчасових положень про пілотування для уникнення блокування. \\
\midrule
\textbf{Залежність від суміжних систем:} Готовність зовнішніх реєстрів (Дія, Мін'юст тощо) до інтеграцій. & Висока & Високий & Проектування архітектури з використанням моків (mocks) на етапі розробки, фіксація SLA взаємодії. \\
\midrule
\textbf{Організаційний супротив:} Слабке залучення команд у процес планування OKR. & Середня & Середній & Систематичне навчання, менторська підтримка, прив'язка OKR до пріоритетів роботи команд. \\
\midrule
\textbf{Технічні збої:} Проблеми стабільності під час перших пілотних запусків у судах. & Середня & Високий & Поетапне розгортання (спочатку 1--2 суди), створення виділеної служби підтримки пілоту, наявність плану відкату (rollback plan). \\
\bottomrule
\end{longtable}

\subsection{Додаток Б: Карта залежностей підсистем (ISO 21502)}
Модернізація підсистем ЄСІКС здійснюється відповідно до таксономії Політик і містить такі ключові залежності (Dependencies):
\begin{enumerate}
    \item \textbf{Фундаментальні сервіси} (Автентифікація КЕП, Рольова модель ABAC/RBAC, Словникова підсистема) є критичною передумовою для старту реалізації \textbf{Господарської діяльності} та \textbf{Судопису}.
    \item \textbf{Інфраструктура документообігу} (PDF-генерація, OCR розпізнавання, QR-кодування) повинна бути завершена до старту пілотування \textbf{Реєстру виконавчих документів та судових рішень}.
    \item \textbf{Інформаційно-пошукова система} (court.gov.ua) та \textbf{Електронний Кабінет} залежать від готовності API-інтерфейсів підсистем \textbf{Судопису} (Справи, Засідання, Авторозподіл).
\end{enumerate}

\subsection{Додаток В: Глосарій термінів та скорочень}
\begin{longtable}{@{}p{3.5cm}p{10.7cm}@{}}
\toprule
\textbf{Термін / Скорочення} & \textbf{Визначення та опис в контексті програми} \\
\midrule
\textbf{ЄСІКС} & Єдина судова інформаційно-телекомунікаційна система. \\
\textbf{eCourt} & Електронний суд (інтерфейсна та функціональна частина ЄСІКС для користувачів). \\
\textbf{OKR} & Objectives and Key Results (Цілі та Ключові Результати) --- методологія стратегічного планування та управління. \\
\textbf{MVP} & Minimum Viable Product (Мінімально життєздатний продукт) --- перша версія підсистеми з достатнім функціоналом для запуску. \\
\textbf{КЕП} & Кваліфікований електронний підпис. \\
\textbf{ДСА} & Державна судова адміністрація України. \\
\textbf{ВРП} & Вища рада правосуддя. \\
\textbf{PMO} & Project Management Office (Офіс управління проектами). \\
\textbf{ABAC / RBAC} & Моделі контролю доступу на основі атрибутів (Attribute-Based) та ролей (Role-Based). \\
\bottomrule
\end{longtable}

\end{document}
